% Các bảng dài tách khỏi thân báo cáo, được \input từ backmatter/phuluc.tex

\section*{Phụ lục F. Ánh xạ vấn đề, hướng giải quyết và mục tiêu}
\label{apdx:anhxa}
\addcontentsline{toc}{section}{Phụ lục F. Ánh xạ vấn đề, hướng giải quyết và mục tiêu}

Bảng~\ref{tab:vande-muctieu} ánh xạ từng vấn đề nêu tại mục~\ref{sec:vande}
tới hướng giải quyết của Farmery và các mục tiêu cụ thể tương ứng tại
mục~\ref{sec:muctieu}.

\begingroup
\small
\renewcommand{\arraystretch}{1.3}
\setlength{\tabcolsep}{5pt}
\begin{longtable}{L{3.3cm} L{7.2cm} L{2.4cm}}
\caption{Ánh xạ giữa vấn đề, hướng giải quyết và mục tiêu của đồ án}
\label{tab:vande-muctieu} \\
\toprule
\rowcolor{tblheadbg}
\textbf{Vấn đề} & \textbf{Hướng giải quyết của Farmery} & \textbf{Mục tiêu} \\
\midrule
\endfirsthead
\toprule
\rowcolor{tblheadbg}
\textbf{Vấn đề} & \textbf{Hướng giải quyết của Farmery} & \textbf{Mục tiêu} \\
\midrule
\endhead
\bottomrule
\endfoot
\textbf{V1.} Cung dồn theo mùa vụ &
Đặt hàng bán sỉ theo hợp đồng và lịch giao để nhà cung cấp chủ động sản xuất; nền tảng chủ động thu mua theo lô ở luồng bán lẻ, tạo thêm một đầu ra ổn định; mô-đun AI dự báo giá theo mùa vụ và vùng trồng. &
MT3, MT4, MT6 \\
\rowcolor{tblaltbg}
\textbf{V2.} Thiếu kênh tiêu thụ &
Hai luồng tiêu thụ song song trên cùng một tập nguồn cung, mở rộng thị trường ra ngoài phạm vi thương lái địa phương; nhà cung cấp không đủ năng lực bán lẻ vẫn tiêu thụ được hàng qua luồng nền tảng tự thu mua. &
MT3, MT4 \\
\textbf{V3.} Hao hụt, thiếu kho lạnh &
Quy trình vận chuyển riêng cho hàng dễ hư hỏng, ưu tiên giao nhanh qua đối tác vận chuyển thứ ba; quản lý tồn kho theo lô với nguyên tắc xuất trước hàng hết hạn trước; giới hạn bán kính giao hàng lẻ để rút ngắn thời gian hàng nằm trên đường. &
MT4, MT5 \\
\rowcolor{tblaltbg}
\textbf{V4.} Thiếu minh bạch, niềm tin &
Xác thực nhà cung cấp và kiểm duyệt chứng nhận; mã QR truy xuất gắn nhật ký canh tác theo lô; hệ thống đánh giá và điểm uy tín; ở luồng bán lẻ, nền tảng trực tiếp chịu trách nhiệm chất lượng thay vì để người mua tự đòi từng người bán nhỏ lẻ. &
MT1, MT2, MT4, MT7 \\
\end{longtable}
\endgroup

\section*{Phụ lục G. Đối chiếu nông sản với các loại hàng hóa khác}
\label{apdx:dacdiem}
\addcontentsline{toc}{section}{Phụ lục G. Đối chiếu nông sản với các loại hàng hóa khác}

Bảng~\ref{tab:dacdiem-hanghoa} đối chiếu nông sản tươi với hàng công nghiệp
và hàng tiêu dùng đóng gói theo từng tiêu chí, kèm ràng buộc thiết kế mà
mỗi khác biệt đặt ra cho hệ thống. Phần phân tích tóm tắt nằm tại
mục~\ref{subsec:sosanhnongsan}.

\begingroup
\small
\renewcommand{\arraystretch}{1.3}
\setlength{\tabcolsep}{5pt}
\begin{longtable}{L{2.6cm} L{3.4cm} L{3.4cm} L{4.6cm}}
\caption{Đối chiếu đặc điểm nông sản với các loại hàng hóa khác}
\label{tab:dacdiem-hanghoa} \\
\toprule
\rowcolor{tblheadbg}
\textbf{Tiêu chí} & \textbf{Nông sản tươi} & \textbf{Hàng công nghiệp, tiêu dùng đóng gói} & \textbf{Ràng buộc đặt ra cho hệ thống} \\
\midrule
\endfirsthead
\toprule
\rowcolor{tblheadbg}
\textbf{Tiêu chí} & \textbf{Nông sản tươi} & \textbf{Hàng công nghiệp, tiêu dùng đóng gói} & \textbf{Ràng buộc đặt ra cho hệ thống} \\
\midrule
\endhead
\bottomrule
\endfoot
Hạn sử dụng & Vài ngày đến vài tuần & Nhiều tháng đến nhiều năm & Quản lý tồn kho theo lô kèm hạn dùng; xuất trước lô hết hạn trước \\
\rowcolor{tblaltbg}
Nguồn cung & Theo mùa vụ, không điều chỉnh được trong ngắn hạn & Theo kế hoạch sản xuất, co giãn theo cầu & Cần dự báo nhu cầu và giá theo mùa vụ, vùng trồng \\
Biến động giá & Mạnh, theo ngày và theo vụ & Ổn định, thay đổi theo chính sách giá & Giá phải cập nhật thường xuyên; biên lợi nhuận dễ bị bào mòn \\
\rowcolor{tblaltbg}
Tính đồng nhất & Khác nhau giữa các lô cùng giống & Đồng nhất theo mã hàng & Dữ liệu mô tả gắn theo lô, không gắn theo mã sản phẩm \\
Đơn vị bán & Khối lượng, bó, thùng; khác nhau giữa lẻ và sỉ & Theo cái, theo hộp & Hỗ trợ nhiều đơn vị tính và quy đổi giữa chúng \\
\rowcolor{tblaltbg}
Khả năng kiểm chứng chất lượng & Người mua không tự kiểm chứng được, kể cả sau khi dùng & Kiểm chứng được bằng quan sát hoặc dùng thử & Bắt buộc có truy xuất nguồn gốc và chứng nhận bên thứ ba \\
Ràng buộc pháp lý & An toàn thực phẩm, kiểm dịch & Chủ yếu về nhãn mác, tiêu chuẩn kỹ thuật & Kiểm duyệt không thể chỉ dựa vào tự động hóa \\
\rowcolor{tblaltbg}
Hao hụt vận chuyển & 20--40\% với rau quả tại Việt Nam & Gần như không đáng kể & Chi phí hao hụt phải tính vào mô hình kinh doanh \\
Điều kiện bảo quản & Phụ thuộc chuỗi lạnh & Nhiệt độ thường & Phụ thuộc năng lực đối tác vận chuyển, nằm ngoài kiểm soát trực tiếp \\
\end{longtable}
\endgroup

\section*{Phụ lục H. Chứng nhận nông sản phổ biến trên thế giới}
\label{apdx:chungnhan}
\addcontentsline{toc}{section}{Phụ lục H. Chứng nhận nông sản phổ biến trên thế giới}

Bảng~\ref{tab:chungnhan} tổng hợp các chứng nhận và giấy tờ mà một lô nông
sản cần có để ra thị trường, kèm mức độ hỗ trợ trong phạm vi đồ án. Phần
phân tích tóm tắt nằm tại mục~\ref{subsec:chungnhan}.

\begingroup
\small
\renewcommand{\arraystretch}{1.3}
\setlength{\tabcolsep}{5pt}
\begin{longtable}{L{3.1cm} L{4.1cm} L{3.6cm} L{3.2cm}}
\caption{Chứng nhận và giấy tờ nông sản theo thị trường}
\label{tab:chungnhan} \\
\toprule
\rowcolor{tblheadbg}
\textbf{Tên} & \textbf{Phạm vi chứng nhận} & \textbf{Thị trường đích} & \textbf{Farmery hỗ trợ} \\
\midrule
\endfirsthead
\toprule
\rowcolor{tblheadbg}
\textbf{Tên} & \textbf{Phạm vi chứng nhận} & \textbf{Thị trường đích} & \textbf{Farmery hỗ trợ} \\
\midrule
\endhead
\bottomrule
\endfoot
\multicolumn{4}{l}{\textbf{\textit{Thực hành sản xuất}}} \\
VietGAP & Thực hành nông nghiệp tốt, khoảng 70 tiêu chí & Nội địa & Có, kiểm duyệt \\
\rowcolor{tblaltbg}
GlobalG.A.P. & Thực hành nông nghiệp tốt, 252 tiêu chuẩn & Xuất khẩu & Có, kiểm duyệt \\
GRASP & Điều kiện lao động, bổ trợ GlobalG.A.P. & Xuất khẩu & Ghi nhận \\
\rowcolor{tblaltbg}
USDA Organic, EU Organic, JAS Organic & Canh tác hữu cơ & Hoa Kỳ, châu Âu, Nhật Bản & Có, kiểm duyệt \\
\multicolumn{4}{l}{\textbf{\textit{An toàn thực phẩm, chế biến và đóng gói}}} \\
HACCP & Phân tích mối nguy và điểm kiểm soát tới hạn & Nội địa và xuất khẩu & Có, kiểm duyệt \\
\rowcolor{tblaltbg}
ISO 22000, FSSC 22000 & Hệ thống quản lý an toàn thực phẩm & Xuất khẩu & Ghi nhận \\
BRCGS Food Safety, IFS Food & Chuẩn do chuỗi bán lẻ lớn yêu cầu & Anh, châu Âu & Ghi nhận \\
\multicolumn{4}{l}{\textbf{\textit{Thị trường ngách}}} \\
\rowcolor{tblaltbg}
Halal, Kosher & Yêu cầu tôn giáo & Thị trường tương ứng & Ghi nhận \\
Fairtrade, Rainforest Alliance & Thương mại công bằng, bền vững & Châu Âu, Bắc Mỹ & Ghi nhận \\
\multicolumn{4}{l}{\textbf{\textit{Giấy tờ đi theo từng lô hàng}}} \\
\rowcolor{tblaltbg}
Kiểm dịch thực vật & Tình trạng dịch hại của lô hàng & Xuất khẩu & Ngoài phạm vi \\
Giấy chứng nhận xuất xứ & Nguồn gốc xuất xứ, ưu đãi thuế quan & Xuất khẩu & Ngoài phạm vi \\
\rowcolor{tblaltbg}
Phiếu kết quả phân tích & Dư lượng thuốc bảo vệ thực vật, kim loại nặng & Nội địa và xuất khẩu & Ngoài phạm vi \\
Mã số vùng trồng, mã cơ sở đóng gói & Định danh nơi trồng và nơi đóng gói; cấp theo Nghị định 38/2026/NĐ-CP & Xuất khẩu chính ngạch & Ghi nhận \\
\rowcolor{tblaltbg}
Điều kiện an toàn thực phẩm, tự công bố sản phẩm & Điều kiện cơ sở và hồ sơ sản phẩm chế biến & Nội địa & Có, bắt buộc với nhóm chế biến sâu \\
OCOP & Xếp hạng sản phẩm đặc trưng địa phương & Nội địa & Có, kiểm duyệt \\
\end{longtable}
\endgroup
